% Filled example -- mirrors docx/abstract-example.docx.
% Build:  latexmk -pdf abstract-example.tex      (runs pdflatex + biber)
\documentclass{confabs}

\title{Microstructure Evolution in Hot-Rolled Steel During Multi-Pass
Deformation}
\author{Anna Kowalska\textsuperscript{1,*}, Jan Nowak\textsuperscript{2},
Maria Lewandowska\textsuperscript{1}}
\affiliation{1}{AGH University of Krakow, Faculty of Metals Engineering,
al.\ Mickiewicza 30, 30-059 Krakow, Poland}
\affiliation{2}{\L{}ukasiewicz Research Network, Institute for Ferrous
Metallurgy, K.\ Miarki 12, 44-100 Gliwice, Poland}
\email{a.kowalska@agh.edu.pl, j.nowak@imz.pl, m.lewandowska@agh.edu.pl}
\keywords{hot rolling, microstructure, recrystallization, finite element
modelling, steel}

\begin{document}
\maketitle

Grain refinement during thermomechanical processing controls the strength and
toughness of structural steels~\cite{humphreys2004}, yet industrial multi-pass
schedules retain strain between passes in a way that single-pass
Johnson--Mehl--Avrami--Kolmogorov (JMAK) calibrations do not
capture~\cite{sellars1979}. This work quantifies how inter-pass time and
temperature govern recrystallization in low-carbon steel. Plane-strain
compression specimens were deformed between \SI{900}{\celsius} and
\SI{1100}{\celsius} following the four-pass schedule of Table~\ref{tab:schedule},
then interrupted-quenched for metallography.

\begin{table}[h]
  \centering
  \caption{Four-pass plane-strain compression schedule.}
  \label{tab:schedule}
  \begin{tabular}{c S[table-format=4.0] S[table-format=1.2] S[table-format=2.0]}
    \toprule
    {Pass} & {Temperature (\si{\celsius})} & {Strain, $\varepsilon$ (--)} &
    {Inter-pass time (s)} \\
    \midrule
    1 & 1100 & 0.30 & 10 \\
    2 & 1050 & 0.30 & 5  \\
    3 & 1000 & 0.25 & 3  \\
    4 & 950  & 0.25 & {--} \\
    \bottomrule
  \end{tabular}
\end{table}

Figure~\ref{fig:kinetics} shows the recrystallized fraction as a function of
inter-pass time at two temperatures. Above \SI{1000}{\celsius} the kinetics are
dominated by inter-pass time rather than peak strain; a JMAK law fitted to
single-pass data over-predicts grain refinement in the multi-pass case by up to
\SI{18}{\percent}~\cite{sellars1979}.

\begin{figure}[h]
  \centering
  \begin{tikzpicture}
    \begin{axis}[
      width=0.86\linewidth, height=6.2cm,
      xlabel={Inter-pass time, $t$ (s)},
      ylabel={Recrystallized fraction, $X$ (\%)},
      xmin=0, xmax=10, ymin=0, ymax=100,
      grid=both, major grid style={gray!25}, minor grid style={gray!10},
      legend pos=north west, legend cell align=left,
      tick label style={font=\footnotesize}, label style={font=\small}]
      \addplot[red, thick, smooth, domain=0:10, samples=120]
        {100*(1-exp(-6.5*(x/10)^1.6))};
      \addlegendentry{\SI{1100}{\celsius}}
      \addplot[blue, thick, smooth, domain=0:10, samples=120]
        {100*(1-exp(-3.2*(x/10)^1.6))};
      \addlegendentry{\SI{1000}{\celsius}}
    \end{axis}
  \end{tikzpicture}
  \caption{Recrystallized fraction $X(t)$ versus inter-pass time at
  \SI{1000}{\celsius} and \SI{1100}{\celsius}.}
  \label{fig:kinetics}
\end{figure}

A retained-strain correction to the JMAK law reconciles laboratory and
industrial bar-mill data, improving grain-size prediction across the
schedule~\cite{kowalska2021}.

\printbibliography[title={References}]
\end{document}
